\chapter{LLD — Calcul et hyperviseurs}
\label{chap:compute}

\begin{todobox}[Chapitre à compléter après identification des CPU et cartes réseau]
Les modèles exacts de CPU et de cartes réseau des R630 ne sont pas encore connus. Ce chapitre présente les décisions déjà prises et les points à compléter.
\end{todobox}

\section{Phase transitoire — Proxmox VE sur R630}
\subsection{Architecture}
\begin{itemize}
  \item Cluster Proxmox VE 8 sur les quatre R630.
  \item Quorum configuré (nœuds impairs ou QDevice sur une VM de management).
  \item Réseau de management sur VLAN 10.
  \item Réseau de VM sur VLAN 20 (bridge vmbr1).
  \item Réseau de migration sur VLAN 40 (bridge vmbr2, dédié).
  \item Stockage local ZFS en miroir (RAIDZ1 ou miroir selon le nombre de disques).
  \item Backups via Proxmox Backup Server (VM dédiée ou serveur historique si réutilisable).
\end{itemize}

\subsection{Rôles des quatre nœuds}
\begin{todobox}[À finaliser après audit]
Selon les CPU et la RAM installés, répartir les rôles :
\begin{itemize}
  \item \texttt{pve01} — nœud principal, héberge la majorité des VM de service.
  \item \texttt{pve02} — nœud secondaire, redondance des VM critiques.
  \item \texttt{pve03} — nœud pour le labo OpenStack (VMs OpenStack en nested virtualization).
  \item \texttt{pve04} — nœud pour les services utilisateurs (pilote).
\end{itemize}
La répartition définitive dépendra des capacités installées.
\end{todobox}

\subsection{Labo OpenStack sur Proxmox}
\begin{itemize}
  \item Déployer trois VM « contrôleurs » (2 vCPU, 8 Go RAM, 40 Go disque).
  \item Déployer deux VM « compute » (4 vCPU, 16 Go RAM, 100 Go disque).
  \item Un réseau de management dédié (VLAN interne Proxmox, non routé).
  \item Un réseau de stockage dédié.
  \item Installer OpenStack via Kolla-Ansible (recommandé) ou DevStack (déconseillé pour cause de non-persistance).
\end{itemize}

\begin{infobox}[Choix Kolla-Ansible]
Kolla-Ansible est retenu pour le labo et pour la production. Il est mature, containerisé, reproductible, et permet de mutualiser les procédures entre les deux phases.
\end{infobox}

\section{Phase cible — OpenStack bare metal}
\subsection{Rôles}
\begin{tabularx}{\textwidth}{@{}lX@{}}
\toprule
\textbf{Rôle} & \textbf{Description} \\
\midrule
Controller & Keystone, Nova API, Neutron server, Glance API, Cinder API, Horizon ($\times 3$ pour HA) \\
Compute    & Nova compute, Neutron agent, Ceph OSD (si convergé) \\
Storage    & Ceph OSD ou nœuds Cinder dédiés selon l'option retenue \\
Network    & Neutron L3 agent si DVR non retenu \\
\bottomrule
\end{tabularx}

\begin{todobox}[À finaliser après audit serveurs historiques]
La répartition dépendra du nombre de serveurs historiques, de leur RAM, de leurs CPU et de leurs disques. La cible peut être totalement ou partiellement convergée.
\end{todobox}